\documentclass[12pt]{article}

\usepackage{amsmath}
\usepackage{german}
\usepackage{graphicx}
\DeclareMathOperator\sgn{sgn}

%\graphicspath{{/home/ritz/Teaching/fcn/fig/}}

%----------------- layout ----------------------------------------

\oddsidemargin=0in
\textwidth=6in
\topmargin=-0.7in
\textheight=9.5in
\baselineskip=16pt
\font\mittel=cmr17
\font\roman=cmr12
\font\klein=cmr9
\font\mgross=cmr10
\newfont{\CAP}{cmcsc10 scaled 900}
\pagestyle{empty}
\parindent=0mm

%----------------- macros ----------------------------------------
\newcommand{\conc}[1]{[{\rm C}]_{\text{#1}}}
\newcommand{\units}[1]{\text{#1}}
\newcommand{\K}{\text{K}}
\newcommand{\Na}{\text{Na}}
\newcommand{\Cl}{\text{Cl}}
\newcommand{\Leak}{\text{leak}}
\newcommand{\water}{\text{H}_2\text{O}}
\newcommand{\tderiv}[1]{\frac{\mathrm{d}#1}{\mathrm{d}t}}
%-----------------------------------------------------------------


\begin{document}

\parbox{2cm}{
\includegraphics[width=1.8cm]{bccnlogo.pdf}
}
\parbox{11cm}{
\begin{center}
\large HUMBOLDT-UNIVERSIT\"AT \hskip 0.1 cm ZU \hskip 0.1 cm BERLIN
\vskip 0.1 true cm
\mgross BERNSTEIN CENTER FOR COMPUTATIONAL NEUROSCIENCE
\end{center}
}
\parbox{2cm}
{
\hfill
\includegraphics[width=1.8cm]{hublogo.pdf}
}

\vskip 0.6 true cm

\leftline{\CAP Humboldt-Universit\"at zu Berlin 
\hfill Phone: 030/2093-9110}
\leftline{\CAP Philippstr. 13 House 6
              \hfill Fax: 030/2093-6771}
\leftline{\CAP \hfill webpage: http://www.bccn-berlin.de/}


\vskip 0.8 true cm
\centerline{\bf Models of Neural Systems, WS 2010/11}
\centerline{\bf Computer Practical 8}
\centerline{Solutions to hand in on: January, 3rd, 2011}


% {\bf Hodgkin-Huxley model}

\vskip 0.8 true cm

\begin{enumerate}

\item {\bf Hodgkin-Huxley model of action potential generation.}
Action potentials (aka spikes) are a dominant feature of nervous systems and
play crucial roles in the computation performed by cortex, thalamus,
cerebellum and associated structures.  The ionic mechanisms of action
potential generation were described for the squid giant axon by Hodgkin
and Huxley. The complete model they proposed is as following:
% Hodgkin and Huxley proposed the following model for action potential
% generation in the squid giant axon:
\begin{equation*}
    C_m \tderiv{V(t)}=-I_{\Na}(V(t), t) - I_{\K}(V(t), t) - I_{\Leak}(V(t),t)
+ I_e
    \label{eq:HHM}
\end{equation*}
where $I_{\Na}(V,t)$ and $I_{\K}(V,t)$ are the sodium and potassium currents
that we already discussed in the previous sheet:
\begin{align*}
    I_{\Na}(V,t) &= \bar{g}_{\Na}m(V,t)^3h(V,t)(V-E_{\Na}), \\
    I_K(V,t) &= \bar{g}_{K} n(V,t)^4 (V-E_K),
    \label{eq:potassium}
\end{align*}
with the reversal potentials $E_{\Na}=50\,\mathrm{mV}$ and
$E_K=-77\,\mathrm{mV}$ and the maximal conductances
$\bar{g}_{\Na}=120\,\mathrm{mS}$ and $\bar{g}_K=36\,\mathrm{mS}$. The gating
variables follow a first-order kinetic equation of the type
\begin{equation*}
    \tderiv{x}=\alpha_x(V)(1-x)-\beta_x(V)x,
    \label{eq:activation}
\end{equation*}
with $x \in \{n,m,h\}$ and the voltage dependent transition rates $\alpha$
and $\beta$:
\begin{align*}
        \alpha_n(V)&=\frac{0.01(V+55)}{1-\exp(-0.1(V+55))}, &
        \beta_n(V)&=0.125\exp(-0.0125(V+65)), \\
        \alpha_m(V)&=0.1\frac{V+40}{1-\exp(-0.1(V+40))}, &
        \beta_m(V)&=4\exp(-0.0556(V+65)), \\
        \alpha_h(V)&=0.07\exp(-0.05(V+65)), &
        \beta_h(V)&=\frac{1}{1+\exp(-0.1(V+35))}.
\end{align*}
$I_{\Leak}(V,t)$ is a passive leakage current due to chloride and other ions:
\begin{equation*}
I_{\Leak}(V,t)=\bar{g}_{\Leak}(V-E_{\Leak})
\end{equation*}
with $\bar{g}_{\Leak}=0.3\,\mathrm{mS}$ and
$E_{\Leak}=-54.387\,\mathrm{mV}$.

\begin{enumerate}
    \item Simulate the Hodgkin-Huxley model by numerically integrating the
        equations for $V$, $n$, $m$ and $h$ (4-dimensional system of ODEs!).
        Take $I_e=0$A, $C_m=1\,\mu\mathrm{F}$ and as initial values take
        $V=-65\,\mathrm{mV}$, $m=0.0529$, $h=0.5961$, and $n=0.3177$. Make
        sure that the integration time step is
        short enough to obtain a stable and accurate solution (test several
        values).
  \item Use a constant external current $I_e=20\,\mu\mathrm{A}$ and plot $V$,
        $I_{\K}$ and $I_{\Na}$ as functions of time. Which currents are
        responsible for different phases of action potential generation?
  \item Plot the firing rate of the model as a function of $I_e$ over
      the range from 0 to 200 $\mu$A. Show that the firing rate
      jumps discontinuously from zero to a finite value when the
      current passes through the minimum value to produce a sustained
      firing. (\textit{Hint:} use a voltage threshold criterion for spike
      detection).
  \item Apply a pulse of a negative current with
      $I_e=-5\,\mu\mathrm{A}$ for 5 ms followed by $I_e=0$ and
      show what happens. Explain the obtained result!
\end{enumerate}

\item {\bf Report.} Submit the plain source code files for all exercises
         and a report containing the source code, all plots and answers to the
         questions asked.

% In the Hodgkin-Huxley formalism, currents flowing through a large number of
% ion
% channels are modelled. The behaviour of single ion channels is more
% complex. Today you model the state of single channels
% using state diagrams and transition rates. In the end you should
% compare the fraction of open channels from your simulations with the
% respective fraction given by the Hodgkin-Huxley model.
% \vskip 0.8 true cm
% 
% \item {\bf Five state potassium channel model.}
% 
% \begin{enumerate}
% \item General Methods
% \begin{enumerate}
%   \item Recapitulate Section 5.7 from the book of Dayan and Abbott. Make sure
%   you understand the state diagrams and the equations describing the
%   state transitions.  \item The transitions rates ($ms^{-1}$) we
%   need are the same we used for the Hodgkin-Huxley model before:
%   \begin{align*}
%       \alpha_n&=0.01\frac{V+55}{1-\exp(-0.1(V+55))}, &
%       \beta_n&=0.125\exp(-0.0125(V+65))\\
%   \end{align*}
% where $V$ is the voltage in mV.
% 
% 
% \end{enumerate}
% \item Assignments
% 
% We assume that the $K^+$ channel consists of four individual subunits,
% which all open and close independently. Thus the channel can be in
% five different states: with $0$, $1$, $2$, $3$, or $4$ subunits being
% open. Current is conducted only in the last state, where all subunits are
% open.
% 
% \begin{enumerate}
% \item Implement a model of $K^+$ channels using the state diagramm
% shown in Figure 5.12. In general, you start with $x$ channels in
% some starting state (e.g. all closed). Then at each time step
% $\Delta t$ you draw for each channel a random number between $0$ and
% $1$ and determine via the state transition rates whether this
% channel stays in its state or changes to another state. The
% probability of changing to a given new state, is its transitions
% rate times $\Delta t$. Make sure that $\Delta t$ is small enough so
% that the probability is always much smaller than $1$!
% 
% \item Test your model by setting an initial membrane voltage and
% ensure that the channel state distributions reaches an equilibrium.
% Then set a new holding potential for some time and observe how the
% channel state distribution changes until the new equilibrium is
% reached. Finally you return to your initial membrane
% potential. Plot the number of open channels as a function of time.
% 
% \item Vary the total number of channels and produce graphs similar to
% the ones in Figure 5.12. Note that the figure shows actual currents, while
% you compare the fraction of open channels. What would you have to
% change to get currents as well?
% \end{enumerate}
% 
% \end{enumerate}
% \bigskip
% \begin{center}
% \includegraphics[scale=.6]{dayan_abbott_fig5_12.pdf}
% \end{center}
% 
% \textbf{Figure 5.12 (from Dayan and Abbott)}. Potassium channel model.
% (Top) State diagram with transition probabilities.  (Bottom) Effect of
% the number of modelled channels on the potassium current.
% 
\end{enumerate}

\vfill
\centerline{\CAP Contact}
\CAP

\begin{tabular}{lll}
Urs Bergmann & Phone: 2093-8924 & Email:
urs.bergmann(at)biologie.hu-berlin.de \\
Matthias Guggenmos & & Email: matthias.guggenmos(at)bccn-berlin.de \\
Richard Kempter \hfill & Phone: 2093-8925 \hfill & Email:
r.kempter(at)biologie.hu-berlin.de \\
\end{tabular}

\end{document}
